\chapter{Regra de vida digital}
\noindent\textit{Vinheta.} Três amigos fazem o mesmo voto: ``\emph{Vou usar menos o celular.}'' Dois fracassam em uma semana. O terceiro cria uma \textbf{regra de vida}: horários, pessoas, práticas e descanso. Em 30 dias, ele \emph{não} virou monge, mas recuperou \textbf{atenção}, \textbf{oração} e \textbf{presença}.

\section{Por que uma regra de vida?}
``Regra'' (do latim \emph{regula}) é um \textbf{ritmo} que apoia a fidelidade no comum da vida. Não é moralismo; é \textbf{mordomia} da atenção diante de telas que disputam nosso amor (Ef~5:15--16).

\begin{nicbox}
\textbf{NIC aplicado}\\
\textbf{Narrativa}: eu sou mais que meu feed; minha história é a de Cristo.\\
\textbf{Identidade}: não sou um perfil; sou filho no Filho.\\
\textbf{Controle}: eu coloco limites nas ferramentas; não entrego o coração a elas.
\end{nicbox}

\section{Princípios simples (4 pilares)}
\begin{itemize}[leftmargin=*,noitemsep]
  \item \textbf{Orações fixas}: breves, em horários chaves (manhã, meio-dia, noite).
  \item \textbf{Janelas de tela}: uso em \textbf{blocos} definidos; notificações fora desses blocos.
  \item \textbf{Comunhão real}: 1 encontro presencial intencional/semana.
  \item \textbf{Descanso sabático}: um período \emph{sem redes} semanal (4--12h), para recuperar o desejo em Deus.
\end{itemize}

\section{Plano de 21 dias (hábitos microscópicos)}
\noindent\textit{Como usar}: faça \textbf{um} bloco por semana (dias 1--7; 8--14; 15--21). Se cair, recomece do ponto da queda, sem culpa.
\subsection*{Dias 1--7 — Atenção que acorda em Deus}
\begin{checklist}
\begin{itemize}[leftmargin=*,noitemsep]
  \item \textbf{Primeiro olhar}: Bíblia antes de tela (Salmo curto em voz alta).
  \item \textbf{Janela matinal}: defina uma hora inicial para abrir apps (ex.: 8h).
  \item \textbf{Café sem tela}: 1 refeição/dia sem celular, com conversa e oração de 30 segundos.
  \item \textbf{Aviso aos seus}: conte a alguém o plano (prestação de contas).
\end{itemize}
\end{checklist}

\subsection*{Dias 8--14 — Ordem no meio do dia}
\begin{checklist}
\begin{itemize}[leftmargin=*,noitemsep]
  \item \textbf{Dois blocos de rede}: ex.: 12h30 e 19h (30 min cada). Fora deles, \emph{apps fechados}.
  \item \textbf{Jejum de encaminhar}: 7 dias sem \emph{forwards} sem checagem 4P (cap.~3).
  \item \textbf{Notificações}: desligue tudo que não seja humano (ligações/agenda).
  \item \textbf{Boa obra pequena}: 1 ato de serviço prático (ligações, visita, oferta).
\end{itemize}
\end{checklist}

\subsection*{Dias 15--21 — Profundidade e sabá digital}
\begin{checklist}
\begin{itemize}[leftmargin=*,noitemsep]
  \item \textbf{Sabá de redes}: escolha 6--12 horas contínuas \emph{sem redes sociais}.
  \item \textbf{Leitura longa}: 30 min de leitura bíblica (ou um Evangelho inteiro em voz alta com alguém).
  \item \textbf{Revisão de apps}: desinstale 1 app que só suga; reorganize a tela inicial para \emph{menos ícones}.
  \item \textbf{Hospitalidade}: convide alguém para uma refeição simples (sem tela na mesa).
\end{itemize}
\end{checklist}

\section{Liturgia doméstica de 15 minutos (modelo)}
\noindent\textit{Sugerido para famílias, casais, amigos ou sozinho. Ajuste livremente.}
\begin{enumerate}[leftmargin=*,nosep]
  \item \textbf{Acolhimento} (1 min): ``O Senhor esteja conosco.'' ``Ele está no meio de nós.''
  \item \textbf{Salmo} (3 min): ler 10--15 versículos em voz alta; uma frase de gratidão.
  \item \textbf{Evangelho} (3 min): um trecho breve (5--10 versículos); silêncio de 30 segundos.
  \item \textbf{Intercessão} (5 min): três pedidos (família, igreja, cidade) + o Pai-Nosso.
  \item \textbf{Envio} (3 min): cada um diz uma \textbf{pequena obediência} para amanhã; bênção (Nm~6:24--26).
\end{enumerate}

\section{Regras de bolso (funcionam na vida real)}
\begin{checklist}
\begin{itemize}[leftmargin=*,noitemsep]
  \item Celular \textbf{fora do quarto}; despertador físico simples.
  \item \textbf{Modo não perturbe} em reuniões, devocionais e refeições.
  \item Um \textbf{dia/semana} sem \emph{shopping} on-line.
  \item \textbf{Sem tela} na primeira e na última meia hora do dia.
\end{itemize}
\end{checklist}

\section{Para líderes: cultura da atenção na igreja}
\begin{checklist}[title={Ações práticas}]
\begin{itemize}[leftmargin=*,noitemsep]
  \item \textbf{Encurtar anúncios}: priorize o culto; dê os detalhes por escrito.
  \item \textbf{Treinar jovens}: oficina de \emph{higiene informacional} (cap.~6) e \emph{checar antes de compartilhar} (cap.~3).
  \item \textbf{Liturgia simples}: leitura pública regular, oração congregacional, silêncio intencional.
  \item \textbf{Semana de sabá digital}: igreja inteira sem redes por 7 dias com devocionais impressos.
\end{itemize}
\end{checklist}

\section{Perguntas para grupos pequenos}
\begin{itemize}[leftmargin=*,noitemsep]
  \item Qual hábito microscópico te daria 80\% do benefício esta semana?
  \item O que te impede de um sabá de redes (6--12h)? O que precisa ser combinado em casa?
  \item Como sua igreja pode apoiar os jovens a recuperarem atenção e oração?
\end{itemize}

\bigskip
\noindent\textit{Próximo capítulo: Famílias e igrejas resilientes --- checklists NIC para tempo de paz e de crise.}
